\documentclass[14pt, a4paper]{extarticle}
\usepackage[T2A]{fontenc}
\usepackage[utf8]{inputenc}
\usepackage[english, russian]{babel}
\usepackage{tempora}
\usepackage{longtable}
\usepackage{array}
\usepackage{ltablex}
\usepackage{ragged2e}
\usepackage{graphicx}
\usepackage{calc}
\usepackage{xcolor} % для цветного текста
\usepackage[left=30mm, right=15mm, top=20mm, bottom=20mm, footskip=12mm]{geometry}

\setlength{\parindent}{0pt}
\setlength{\parskip}{6pt plus 2pt minus 1pt}

\newcolumntype{P}[1]{>{\RaggedRight\arraybackslash}p{#1 - 2\tabcolsep - 2\arrayrulewidth}}
\newcolumntype{C}[1]{>{\centering\arraybackslash}p{#1 - 2\tabcolsep - 2\arrayrulewidth}}

\begin{document}

\pagestyle{empty}

% ------------------------- Заголовок документа -------------------------

\begin{center}
    {\small МИНИСТЕРСТВО НАУКИ И ВЫСШЕГО ОБРАЗОВАНИЯ РОССИЙСКОЙ ФЕДЕРАЦИИ}

    \resizebox{\textwidth}{!}{\fontsize{8pt}{10pt}\selectfont ФЕДЕРАЛЬНОЕ
        ГОСУДАРСТВЕННОЕ АВТОНОМНОЕ ОБРАЗОВАТЕЛЬНОЕ УЧРЕЖДЕНИЕ ВЫСШЕГО ОБРАЗОВАНИЯ}

    {\small \textbf{«МОСКОВСКИЙ ПОЛИТЕХНИЧЕСКИЙ УНИВЕРСИТЕТ»}} \\
    {\small \textbf{(МОСКОВСКИЙ ПОЛИТЕХ)}}
\end{center}

\begin{tabularx}{\textwidth}{@{} p{0.53\textwidth} p{0.43\textwidth} @{}}
    &
    \raggedright
    УТВЕРЖДАЮ                          \\
    заведующая кафедрой                \\
    «Инфокогнитивные технологии»       \\[0.3cm]

    \rule{3cm}{0.4pt} / Е.~А.~Пухова / \\[0.1cm]

    «\rule{0.8cm}{0.4pt}» \rule{3cm}{0.4pt} 2026~г.
\end{tabularx}

% -------------------------- Общая информация --------------------------

\begin{center}
    { \large \textbf{Разработка веб-приложения для управления очередями студентов (на сдачу лабораторных работ и посещений)}}

    \textbf{на выполнение курсовой работы (проекта)}

    обучающимся группы 241-326: \\
    Дудченко Алёне Алексеевне и Факову Станиславу Юрьевичу,

    направления подготовки 09.03.01 «Информатика и вычислительная техника»

    по дисциплине «Разработка веб-приложений»

    на тему «Разработка веб-приложения для управления очередями студентов (на сдачу лабораторных работ и посещений)»
\end{center}

1. Исходные данные к работе (проекту): информационные ресурсы в сети интернет, научные публикации в открытой печати.

2. Содержание задания по курсовой работе (проекту) -- перечень вопросов, подлежащих разработке:

% ------------------------------- Задачи -------------------------------
{
\renewcommand{\arraystretch}{1.2}
\small
\begin{longtable}{|P{0.55\textwidth}|C{0.2\textwidth}|C{0.25\textwidth}|}
    \hline
    \textbf{Разрабатываемый вопрос} & \textbf{Срок} & \textbf{Исполнитель} \\
    \hline
    \endhead
    \hline
    \endfoot

    \multicolumn{3}{|P{\textwidth}|}{\textbf{Раздел 1. Анализ предметной области}} \\ \hline
        1.1 Анализ существующих инструментов организации учебных очередей & 30.03.2026 & Алёна и Станислав \\ \hline
        \hspace*{1em}1.1.1 Анализ существующих алгоритмов организации учебных очередей и обоснование эффективности разработки & 30.03.2026 & \textcolor{blue}{Станислав} \\ \hline
        \hspace*{1em}1.1.2 Анализ существующих инструментов организации учебных очередей и обоснование эффективности разработки & 30.03.2026 & \textcolor{red}{Алёна} \\ \hline
        1.2 Обоснование выбора стека технологий & 06.04.2026 & Алёна и Станислав \\ \hline
        \hspace*{1em}1.2.1 Выбор клиентских технологий (HTML5, CSS3, Jinja2, JS, инструменты вёрстки) & 06.04.2026 & \textcolor{red}{Алёна} \\ \hline
        \hspace*{1em}1.2.2 Выбор серверных технологий (Python, Flask, SQLAlchemy, PostgreSQL) & 06.04.2026 & \textcolor{blue}{Станислав} \\ \hline
        1.3 Формулировка целей, задач и научной новизны (совместно) & 13.04.2026 & Алёна и Станислав \\ \hline
        \hspace*{1em}1.3.1 Описание архитектурных решений и требований к БД & 13.04.2026 & \textcolor{blue}{Станислав} \\ \hline
        \hspace*{1em}1.3.2 Описание новизны алгоритма очереди и метрик эффективности разработки & 13.04.2026 & \textcolor{red}{Алёна} \\ \hline

    \multicolumn{3}{|P{\textwidth}|}{\textbf{Раздел 2. Проектирование}} \\ \hline
        2.1 Анализ целевой аудитории и User Stories & 20.04.2026 & Алёна и Станислав \\ \hline
        \hspace*{1em}2.1.1 Теоретический анализ целевой аудитории и разработка User Stories & 20.04.2026 & \textcolor{blue}{Станислав} \\ \hline
        \hspace*{1em}2.1.2 Описание ролей и сценариев взаимодействия Студент-Преподаватель-Администратор & 20.04.2026 & \textcolor{red}{Алёна} \\ \hline
        2.2 Разработка UML Use Case & 27.04.2026 & Алёна и Станислав \\ \hline
        \hspace*{1em}2.2.1 Для роли «Студент» & 27.04.2026 & \textcolor{blue}{Станислав} \\ \hline
        \hspace*{1em}2.2.2 Для роли «Администратор и Преподаватель» & 27.04.2026 & \textcolor{red}{Алёна} \\ \hline
        2.3 Проектирование модели данных & 04.05.2026 & Алёна и Станислав \\ \hline
        \hspace*{1em}2.3.1 Проектирование ER-диаграммы для сущностей «Администратор и Преподаватель», логической схемы БД & 04.05.2026 & \textcolor{red}{Алёна} \\ \hline
        \hspace*{1em}2.3.2 Проектирование ER-диаграммы для сущности «Студент», физической схемы БД & 04.05.2026 & \textcolor{blue}{Станислав} \\ \hline
        2.4 Разработка макетов Wireframes & 11.05.2026 & Алёна и Станислав \\ \hline
        \hspace*{1em}2.4.1 Разработка макета Wireframe для очереди и Студента & 11.05.2026 & \textcolor{blue}{Станислав} \\ \hline
        \hspace*{1em}2.4.2 Разработка макета Wireframe для Преподавателя и администратора & 11.05.2026 & \textcolor{red}{Алёна} \\ \hline

    \multicolumn{3}{|P{\textwidth}|}{\textbf{Раздел 3. Разработка}} \\ \hline
        3.1 Разработка базовой структуры приложения и вёрстка шаблонов страниц & 18.05.2026 & Алёна и Станислав \\ \hline
        \hspace*{1em}3.1.1 Разработка шаблона страницы преподавателя и администратора, написание логики связи "клиент-сервер" & 18.05.2026 & \textcolor{red}{Алёна} \\ \hline
        \hspace*{1em}3.1.2 Разработка шаблона страницы студента и написание логики очереди & 18.05.2026 & \textcolor{blue}{Станислав} \\ \hline
        3.2 Реализация модуля авторизации и разграничения прав доступа & 25.05.2026 & \\ \hline
        \hspace*{1em}3.2.1 Реализация модуля авторизации для Преподавателя и Студента & 25.05.2026 & \textcolor{blue}{Станислав} \\ \hline
        \hspace*{1em}3.2.2 Реализация модуля авторизации для администратора, разграничения прав доступа для всех ролей & 25.05.2026 & \textcolor{red}{Алёна} \\ \hline
        3.3 Разработка CRUD интерфейса & 08.06.2026 & Алёна и Станислав \\ \hline
        \hspace*{1em}3.3.1 Разработка CRUD взаимодействия "Клиент-Сервер" для очередей & 08.06.2026 & \textcolor{red}{Алёна} \\ \hline
        \hspace*{1em}3.3.2 Разработка CRUD взаимодействия "Клиент-Сервер" для сущностей & 08.06.2026 & \textcolor{blue}{Станислав} \\ \hline
        3.4 Загрузка работ и выгрузка отчётов & 09.06.2026 & \textcolor{blue}{Станислав} \\ \hline
        3.5 Вёрстка адаптивного интерфейса (Jinja2, HTML, CSS) & 10.06.2026 & \textcolor{red}{Алёна} \\ \hline

    \multicolumn{3}{|P{\textwidth}|}{\textbf{Раздел 4. Оформление}} \\ \hline
        4.1 Презентация Git-репозитория & 11.06.2026 & Алёна и Станислав \\ \hline
        4.2 Деплой на облачный хостинг & 12.06.2026 & Алёна и Станислав \\ \hline
        4.3 Оформление отчета о проделанной работе & 13.06.2026 & Алёна и Станислав \\ \hline
\end{longtable}
}

% -------------- Сведения о выдаче и принятии задания -------------------

Руководитель курсовой работы (проекта): \\
старший преподаватель кафедры «Инфокогнитивные технологии»

{
\renewcommand{\arraystretch}{2}
\noindent
\begin{tabularx}{\textwidth}{@{} P{0.45\textwidth} P{0.3\textwidth} P{0.25\textwidth} @{}}
    «\rule{0.8cm}{0.4pt}» \quad \rule{2cm}{0.4pt} \quad 2026~г. &
    \centering \rule{3.5cm}{0.4pt} \par \footnotesize (подпись) &
    \hfill А.~С.~Кружалов \\
\end{tabularx}

\begin{tabularx}{\textwidth}{@{} P{0.45\textwidth} P{0.55\textwidth} @{}}
    Дата выдачи задания           & 
    \hfill «\underline{30}» \quad \underline{марта} \quad 2026~г. \\
    Дата сдачи выполненной работы & 
    \hfill «\underline{13}» \quad \underline{июня} \quad 2026~г. \\
\end{tabularx}
}

Задание приняли к исполнению

{
\renewcommand{\arraystretch}{1.5}
\begin{tabularx}{\textwidth}{@{} P{0.45\textwidth} P{0.3\textwidth} P{0.25\textwidth} @{}}
    «\rule{0.8cm}{0.4pt}» \quad \rule{2cm}{0.4pt} \quad 2026~г. &
    \centering \rule{3.5cm}{0.4pt} \par \footnotesize (подпись) &
    \hfill А. А. Дудченко \\
    «\rule{0.8cm}{0.4pt}» \quad \rule{2cm}{0.4pt} \quad 2026~г. &
    \centering \rule{3.5cm}{0.4pt} \par \footnotesize (подпись) &
    \hfill С. Факов \\
\end{tabularx}
}

\end{document}
